% Part: history
% Chapter: biographies
% Section: kurt-goedel

\documentclass[../../../include/open-logic-section]{subfiles}

\begin{document}

\olfileid{his}{bio}{god}

\olsection{كورت غودل}

\olphoto{goedel-kurt}{Kurt G{"o}del}

وُلد كورت غودل (\textsc{ger}-dle) في 28 أبريل 1906 في برون في
الإمبراطورية النمساوية المجرية (وهي الآن برنو في جمهورية التشيك). وبسبب
طبيعته الفضولية وذكائه، كثيراً ما كانت أسرته تدعو كورت الصغير «Der kleine
Herr Warum» (السيد «لماذا» الصغير). وقد تفوق في دراسته منذ المدرسة
الابتدائية، ولم يحصل على درجة أقل من الدرجة القصوى إلا في الرياضيات. وكان
غودل يتغيب عن المدرسة كثيراً بسبب اعتلال صحته، وأُعفي من التربية البدنية.
وشُخّص في طفولته بالحمى الروماتيزمية. وظل طوال حياته يعتقد أنها أضرت قلبه
إضراراً دائماً، على الرغم من أن التقييم الطبي قال خلاف ذلك.

بدأ غودل الدراسة في جامعة فيينا عام 1924، وأتم دراسته للدكتوراه عام~1929.
وكان ينوي أولاً دراسة الفيزياء، لكن اهتمامه تحول سريعاً إلى الرياضيات،
ولا سيما المنطق، ويرجع ذلك جزئياً إلى تأثير الفيلسوف رودلف كارناب. وقد
برهنت أطروحته، التي كتبها بإشراف هانس هان، مبرهنة تمام منطق المحمولات من
الرتبة الأولى مع الهوية \citep{Godel1929}. وبعد عام واحد فقط حصل على أشهر
نتائجه: مبرهنتي عدم التمام الأولى والثانية (المنشورتين في
\citealt{Godel1931}). وكان غودل في أثناء إقامته في فيينا وثيق الصلة بحلقة
فيينا، وهي جماعة من الفلاسفة ذوي التوجه العلمي كان من بينهم كارناب، وقد
تأثر عملهم تأثراً خاصاً بنتائج غودل.

تزوج غودل عام 1938 أديله نيمبورسكي. ولم يرضَ والداه؛ فهي لم تكن تكبره بست
سنوات ومطلقةً فحسب، بل كانت تعمل راقصةً في ملهى ليلي. غير أن الضغوط
الاجتماعية لم تؤثر في غودل، وظلا سعيدين في زواجهما حتى وفاته.

بعد أن ضمت ألمانيا النازية النمسا عام 1938، هاجر غودل وأديله إلى الولايات
المتحدة، حيث تولى منصباً في معهد الدراسات المتقدمة في برينستون بولاية
نيوجيرسي. وعلى الرغم من انطوائه وغرابة أطواره، كانت سنواته في برينستون
حافلةً بالتعاون والإنتاج. فنشر مقالات في نظرية المجموعات والفلسفة
والفيزياء. ومن أبرز ما حدث أنه أقام صداقة متينة على نحو خاص مع زميله في
المعهد ألبرت أينشتاين.

تدهورت صحة غودل النفسية في سنواته الأخيرة. وحين دخلت زوجته المستشفى عام
1977، لم تعد قادرةً على إعداد طعامه. وقد عانى مشكلات نفسية طوال حياته، ثم
استسلم للارتياب المرضي. وإذ كان يخشى خوفاً قاتلاً أن يُسمَّم، امتنع عن
الأكل. وتوفي جوعاً في 14 يناير 1978 في برينستون.

\begin{reading}
للاطلاع على سيرة كاملة لحياة غودل، انظر \citet{Dawson1997}. ولمزيد من
المواد المتعلقة بسيرته، فضلاً عن مقالات في إسهاماته في المنطق والفلسفة،
انظر \citet{Wang1990} و\citet{Baaz2011} و\citet{Takeuti2003} و
\citet{Sigmund2007}.

أطروحة غودل للدكتوراه متاحة بالألمانية الأصلية \citep{Godel1929}.
والنص الأصلي لمبرهنتي عدم التمام هو \citep{Godel1931}. وجميع كتابات غودل
المنشورة وغير المنشورة، مع مختارات من مراسلاته، متاحة بالإنجليزية في
\emph{أوراقه المجموعة} \citet{Godel1986,Godel1990}.

للمعالجة المفصلة لمبرهنتي غودل في عدم التمام، انظر \citet{Smith2013}.
ولمناقشة فلسفية غير رسمية لمبرهنتيه، انظر تدوينة مارك لينسنماير الصوتية
\citep{Linsenmayer2014}.
\end{reading}

\end{document}
